% ============================================================================
% LANGENTWURF LE-03c: Los colores + Adjektivangleichung
% Spanisch Klasse 7 - Sequenz 3: Mi casa
% ============================================================================
% Tier: 1 (vollständig)
% Doppelstunde: 3c (Std. 5-6 von 8)
% Fachfarbe: #c41e3a (Karminrot)
% ============================================================================

\documentclass[11pt,a4paper]{article}

\usepackage[utf8]{inputenc}
\usepackage[T1]{fontenc}
\usepackage[ngerman]{babel}
\usepackage{geometry}
\usepackage{fancyhdr}
\usepackage{lastpage}
\usepackage{xcolor}
\usepackage{tcolorbox}
\usepackage{enumitem}
\usepackage{tabularx}
\usepackage{longtable}
\usepackage{array}
\usepackage{booktabs}
\usepackage{amssymb}

% ============================================================================
% KONFIGURATION
% ============================================================================

\newcommand{\fach}{Spanisch}
\newcommand{\fachfarbe}{c41e3a}
\newcommand{\jahrgangsstufe}{7}
\newcommand{\schulart}{Gesamtschule}
\newcommand{\stundenthema}{Los colores + Adjektivangleichung}
\newcommand{\reihenthema}{Mi casa}
\newcommand{\stundennummer}{3c}
\newcommand{\reihenstunden}{4}
\newcommand{\gession}{A1}
\newcommand{\lehrwerk}{Qué pasa 1}
\newcommand{\lehrwerkseite}{S. 88-90}
\newcommand{\lerngruppengroesse}{24}

% ============================================================================
% LAYOUT
% ============================================================================
\geometry{left=2.5cm, right=2.5cm, top=2.5cm, bottom=2.5cm}
\definecolor{fach}{HTML}{\fachfarbe}
\definecolor{gray}{HTML}{666666}
\definecolor{lightgray}{HTML}{f5f5f5}
\definecolor{successgreen}{HTML}{2a7a4b}
\definecolor{warningorange}{HTML}{b35c00}
\definecolor{basis}{HTML}{2a7a4b}
\definecolor{standard}{HTML}{2c5aa0}
\definecolor{erweiterung}{HTML}{7b1fa2}

\tcbuselibrary{skins,breakable}

\newtcolorbox{infobox}[1][]{
    colback=lightgray,
    colframe=fach,
    fonttitle=\bfseries,
    title=#1,
    breakable
}

\newtcolorbox{scorebox}{
    colback=white,
    colframe=fach,
    fonttitle=\bfseries,
    title=Didaktik-Score,
    breakable
}

\pagestyle{fancy}
\fancyhf{}
\fancyhead[L]{\small\textcolor{gray}{\fach{} -- Jg. \jahrgangsstufe{} -- \stundenthema}}
\fancyhead[R]{\small\textcolor{gray}{LE-\stundennummer{} | Tier 1}}
\fancyfoot[C]{\small\thepage{} / \pageref{LastPage}}
\renewcommand{\headrulewidth}{0.4pt}

% Differenzierungsmarker
\newcommand{\B}{\textcolor{basis}{\textbf{[B]}}}
\newcommand{\Std}{\textcolor{standard}{\textbf{[S]}}}
\newcommand{\E}{\textcolor{erweiterung}{\textbf{[E]}}}

% ============================================================================
% DOKUMENT
% ============================================================================
\begin{document}

% ============================================================================
% DECKBLATT
% ============================================================================
\begin{titlepage}
    \centering
    \vspace*{2cm}
    {\large\textcolor{gray}{\schulart}}\\[2cm]
    {\Huge\textcolor{fach}{\textbf{Langentwurf}}}\\[0.5cm]
    {\large LE-\stundennummer{} | Tier 1 (vollständig)}\\[2cm]
    {\LARGE\textbf{\stundenthema}}\\[0.5cm]
    {\large Sequenz: \reihenthema}\\[2cm]
    \begin{tabular}{rl}
        \textbf{Fach:} & \fach{} (A1) \\[0.3cm]
        \textbf{Klasse:} & \jahrgangsstufe \\[0.3cm]
        \textbf{Lehrwerk:} & \lehrwerk, \lehrwerkseite \\[0.3cm]
        \textbf{Doppelstunde:} & \stundennummer{} (Std. 5-6 von 8) \\[0.3cm]
    \end{tabular}
    \vfill
\end{titlepage}

\tableofcontents
\newpage

% ============================================================================
% 1. BEDINGUNGSANALYSE
% ============================================================================
\section{Bedingungsanalyse}

\subsection{Lerngruppe}
\begin{tabular}{ll}
    Kursgröße: & \lerngruppengroesse{} SuS \\
    Jahrgangsstufe: & \jahrgangsstufe \\
    Sprachniveau: & A1 (GER) -- Anfänger \\
    Fremdsprache: & 2. Fremdsprache (nach Englisch) \\
\end{tabular}

\subsection{Vorwissen aus 03a/03b}
\begin{itemize}
    \item Zimmer und Räume im Haus: \textit{la cocina, el salón, el dormitorio, el baño, el jardín}
    \item Möbel und Einrichtungsgegenstände: \textit{la mesa, la silla, el sofá, la cama, el armario}
    \item Präpositionen: \textit{en, sobre, debajo de, al lado de, delante de, detrás de}
    \item Verb \textit{estar} für Ortsangaben
    \item Verb \textit{hay} (es gibt)
\end{itemize}

\subsection{Differenzierung (intern)}
\begin{tabular}{|l|p{3.5cm}|p{3.5cm}|p{3.5cm}|}
\hline
& \B{} \textbf{Basis} & \Std{} \textbf{Standard} & \E{} \textbf{Erweiterung} \\
\hline
\textbf{Ziel} & BR: Rezeption & MR: Produktion mit Hilfe & GY: Freie Produktion \\
\hline
\textbf{Farben} & Zuordnung Bild-Wort & Farben selbst nennen & Nuancen unterscheiden \\
\hline
\textbf{Adj.-Angleichung} & Einsetzen mit Tabelle & Selbstständig angleichen & Sätze frei formulieren \\
\hline
\end{tabular}

% ============================================================================
% 2. SACHANALYSE
% ============================================================================
\section{Sachanalyse}

\subsection{Sprachliche Strukturen}

\textbf{Farbadjektive:}
\begin{center}
\begin{tabular}{|l|l|l|}
\hline
\textbf{Spanisch} & \textbf{Deutsch} & \textbf{Typ} \\
\hline
rojo/a & rot & veränderlich \\
\hline
azul & blau & unveränderlich \\
\hline
verde & grün & unveränderlich \\
\hline
amarillo/a & gelb & veränderlich \\
\hline
negro/a & schwarz & veränderlich \\
\hline
blanco/a & weiß & veränderlich \\
\hline
marrón & braun & unveränderlich \\
\hline
gris & grau & unveränderlich \\
\hline
rosa & rosa & unveränderlich \\
\hline
naranja & orange & unveränderlich \\
\hline
\end{tabular}
\end{center}

\textbf{Adjektivangleichung:}
\begin{itemize}
    \item Adjektive auf \textbf{-o} ändern sich nach Genus und Numerus
    \item Adjektive auf \textbf{-e} oder Konsonant ändern nur den Numerus
    \item Farbadjektive stehen \textbf{nach} dem Nomen
\end{itemize}

\begin{center}
\begin{tabular}{|l|c|c|}
\hline
& \textbf{Singular} & \textbf{Plural} \\
\hline
\textbf{maskulin} & el sofá \textbf{rojo} & los sofás \textbf{rojos} \\
\hline
\textbf{feminin} & la silla \textbf{roja} & las sillas \textbf{rojas} \\
\hline
\end{tabular}
\end{center}

\subsection{Kontrastive Betrachtung}
\begin{itemize}
    \item \textbf{Positiv:} Konzept der Adjektivangleichung aus Deutsch/Englisch bekannt
    \item \textbf{Interferenz:} Im Deutschen stehen Adjektive \textit{vor} dem Nomen, im Spanischen meist \textit{dahinter}
    \item \textbf{Interferenz:} Deutsche SuS vergessen oft die Angleichung (*el sofá roja)
    \item \textbf{Falsche Freunde:} \textit{negro} bedeutet schwarz, nicht negativ
\end{itemize}

\subsection{Kulturelle Aspekte}
\begin{itemize}
    \item Farbsymbolik: Rot in Spanien = Leidenschaft, Stierkampf
    \item Typische Hausfarben in Spanien: Weiß (Andalusien), Terrakotta
    \item Farbbezeichnungen variieren regional (azul marino, azul celeste)
\end{itemize}

% ============================================================================
% 3. DIDAKTISCHE ANALYSE
% ============================================================================
\section{Didaktische Analyse}

\subsection{Stellung in der Sequenz}
Dies ist die \textbf{3. Doppelstunde} (Std. 5-6 von 8) der Sequenz ``\reihenthema''.

\textbf{Sequenzaufbau:}
\begin{itemize}
    \item 3a: Zimmer im Haus (\textit{las habitaciones})
    \item 3b: Möbel + Präpositionen (\textit{los muebles})
    \item \textbf{3c: Farben + Adjektivangleichung (diese Stunde)}
    \item 3d: Mi casa ideal + Sequenztest
\end{itemize}

\subsection{Kompetenzziele (Rahmenplan MV)}
\begin{itemize}
    \item \textbf{Sprachlich:} Adjektive an Genus und Numerus angleichen
    \item \textbf{Kommunikativ:} Gegenstände mit Farben beschreiben (Sprechen, Schreiben)
    \item \textbf{Interkulturell:} Wohnkultur und Farbgestaltung in Spanien
    \item \textbf{Methodisch:} Wortfeldarbeit (systematische Vokabelerschließung)
\end{itemize}

\subsection{Legitimation}
\textbf{Rahmenplan Spanisch MV:}
\begin{itemize}
    \item Verfügen über sprachliche Mittel: Grammatik -- Adjektivangleichung
    \item Schreiben: Gegenstände beschreiben
    \item Sprechen: Über die eigene Wohnung sprechen
\end{itemize}

% ============================================================================
% 4. STUNDENZIELE
% ============================================================================
\section{Stundenziele}

\subsection{Grobziel}
Die SuS erweitern ihre \textbf{sprachliche Kompetenz}, indem sie die 10 wichtigsten Farbadjektive kennenlernen und die Regeln der \textbf{Adjektivangleichung} anwenden, um Gegenstände in ihrem Zimmer zu beschreiben.

\subsection{Feinziele (operationalisiert)}
\begin{tabular}{|c|p{8cm}|c|c|}
\hline
\textbf{Nr.} & \textbf{Die SuS können...} & \textbf{AFB} & \textbf{Diff.} \\
\hline
FZ1 & die 10 Farbadjektive (rojo, azul, verde, amarillo, negro, blanco, marrón, gris, rosa, naranja) benennen und aussprechen & I & alle \\
\hline
FZ2 & Farbadjektive den Nomen korrekt nachstellen & I & alle \\
\hline
FZ3 & veränderliche Adjektive (-o/-a) an Genus und Numerus des Nomens angleichen & II & alle \\
\hline
FZ4 & Gegenstände im Zimmer mit Farben beschreiben (``El sofá es rojo'', ``Las sillas son azules'') & II & \Std{}/\E{} \\
\hline
FZ5 & ihr Traumzimmer (mi habitación ideal) mit Farben beschreiben & III & \E{} \\
\hline
\end{tabular}

% ============================================================================
% 5. METHODISCHE ANALYSE
% ============================================================================
\section{Methodische Analyse}

\subsection{Medien}
\begin{itemize}
    \item Tafel/Whiteboard
    \item Lehrwerk: \lehrwerk, \lehrwerkseite
    \item Arbeitsblatt: AB-03c.pdf
    \item Farbkarten (10 Farben, laminiert)
    \item Bildkarten: Möbel in verschiedenen Farben
    \item Optional: LP-03c.html (digital)
\end{itemize}

\subsection{Sozialformen}
\begin{itemize}
    \item Plenum: Einführung Farben, Regelerarbeitung
    \item Partnerarbeit: Farben-Memory, Beschreibungsspiel
    \item Einzelarbeit: AB-Aufgaben, ``Mi habitación''
\end{itemize}

% ============================================================================
% 6. VERLAUFSPLANUNG
% ============================================================================
\section{Verlaufsplanung}

\textbf{1. Stunde (45 Min.) -- Schwerpunkt: Farben lernen}

\begin{longtable}{|p{1cm}|p{1.5cm}|p{4cm}|p{3cm}|p{2cm}|p{2cm}|}
\hline
\textbf{Zeit} & \textbf{Phase} & \textbf{Lehrerhandeln} & \textbf{Schülerhandeln} & \textbf{SF} & \textbf{Medien} \\
\hline
\endhead

0--5 & Einstieg & Farbige Gegenstände zeigen: ``¿De qué color es?'' & Vermutungen äußern & Plenum & Objekte \\
\hline

5--15 & Input & Farben einführen mit Karten, Tafelanschrieb & Nachsprechen, chorisch & Plenum & Karten \\
\hline

15--25 & Übung 1 & Farben-Drill: Karten zeigen, Farbe abfragen & Schnell antworten: ``¡Rojo!'' & Plenum & Karten \\
\hline

25--35 & Übung 2 & AB-03c Merkkasten erklären, Farben eintragen & Ausfüllen aus Lernpfad/Tafel & EA & AB \\
\hline

35--40 & Sicherung & Farben an Bildern abfragen & Antworten: ``La mesa es marrón'' & Plenum & Bilder \\
\hline

40--45 & Überleitung & Adjektivangleichung ankündigen: ``rojo vs. roja'' & Aufmerksam zuhören & Plenum & Tafel \\
\hline
\end{longtable}

\vspace{1em}

\textbf{2. Stunde (45 Min.) -- Schwerpunkt: Adjektivangleichung}

\begin{longtable}{|p{1cm}|p{1.5cm}|p{4cm}|p{3cm}|p{2cm}|p{2cm}|}
\hline
\textbf{Zeit} & \textbf{Phase} & \textbf{Lehrerhandeln} & \textbf{Schülerhandeln} & \textbf{SF} & \textbf{Medien} \\
\hline
\endhead

0--5 & Aktivierung & Farben-Quiz: ``¿De qué color es el sofá?'' & Schnell antworten & Plenum & Bilder \\
\hline

5--15 & Input & Adjektivangleichung erklären: -o/-a, Plural, Tabelle & Mitschreiben, Beispiele bilden & Plenum & Tafel \\
\hline

15--25 & Übung 1 & AB-03c Aufgabe 2: Adjektive angleichen & Einsetzen: rojo $\rightarrow$ roja & EA & AB \\
\hline

25--35 & Übung 2 & PA: Zimmer beschreiben ``En mi habitación hay...'' & Partner beschreibt, anderer rät & PA & -- \\
\hline

35--40 & Produktion & AB-03c Aufgabe 4: ``Mi casa ideal'' & Schreiben mit Farben & EA & AB \\
\hline

40--45 & Abschluss & Zusammenfassung, HA erklären & Aufschreiben, vorlesen & Plenum & -- \\
\hline
\end{longtable}

\textbf{Legende:} EA = Einzelarbeit, PA = Partnerarbeit, SF = Sozialform

% ============================================================================
% 7. ERWARTETE SCHWIERIGKEITEN
% ============================================================================
\section{Erwartete Schwierigkeiten}

\begin{tabular}{|p{5cm}|p{8cm}|}
\hline
\textbf{Schwierigkeit} & \textbf{Reaktion} \\
\hline
Adjektiv vor Nomen (*rojo sofá) & Bewusstmachung: Spanisch = Adjektiv NACH Nomen; Drill mit Beispielen \\
\hline
Fehlende Angleichung (*la silla rojo) & Tabelle als Hilfe, Merkspruch: ``Nomen bestimmt, Adjektiv folgt'' \\
\hline
Verwechslung marrón/naranja & Beide unveränderlich (keine Angleichung), Eselsbrücke: ``Marrón bleibt marrón'' \\
\hline
Aussprache rr in marrón & Drill: rollender R-Laut, Kontrastübung: pero vs. perro \\
\hline
Plural der unveränderlichen Adj. & Regel: nur -s/-es anhängen (verde $\rightarrow$ verdes, marrón $\rightarrow$ marrones) \\
\hline
\end{tabular}

% ============================================================================
% 8. HAUSAUFGABE
% ============================================================================
\section{Hausaufgabe}

\begin{itemize}
    \item Lehrwerk S. 89, Übung 5 (Farben zuordnen)
    \item Vokabeln lernen: 10 Farben + 5 Möbelstücke
    \item Optional: LP-03c.html bearbeiten (Adjektivangleichung)
\end{itemize}

% ============================================================================
% 9. LÖSUNGEN
% ============================================================================
\section{Lösungen AB-03c}

\textbf{Merkkasten 1 (Farben):}
\begin{itemize}
    \item rojo = rot
    \item azul = blau
    \item verde = grün
    \item amarillo = gelb
    \item negro = schwarz
    \item blanco = weiß
    \item marrón = braun
    \item gris = grau
    \item rosa = rosa
    \item naranja = orange
\end{itemize}

\textbf{Merkkasten 2 (Adjektivangleichung):}
\begin{itemize}
    \item Adjektive stehen NACH dem Nomen
    \item Adjektive auf -o ändern Endung: -o/-a/-os/-as
    \item Andere Adjektive: nur Plural (-s/-es)
\end{itemize}

\textbf{Aufgabe 1 (Farben zuordnen):} (4P)
\begin{itemize}
    \item rojo = rot (C)
    \item azul = blau (A)
    \item verde = grün (D)
    \item amarillo = gelb (B)
\end{itemize}

\textbf{Aufgabe 2 (Adjektiv anpassen):} (4P)
\begin{enumerate}[label=\alph*)]
    \item la mesa \textbf{blanca}
    \item los sofás \textbf{rojos}
    \item las sillas \textbf{azules}
    \item el armario \textbf{negro}
\end{enumerate}

\textbf{Aufgabe 3 (Zimmer beschreiben):} (6P)
\begin{itemize}
    \item En mi habitación hay una cama \textbf{blanca}.
    \item El armario es \textbf{marrón}.
    \item Las cortinas son \textbf{verdes}.
    \item (weitere individuelle Lösungen möglich)
\end{itemize}

\textbf{Aufgabe 4 (Mi casa ideal):} (6P)

Bewertungskriterien:
\begin{itemize}
    \item Mindestens 3 Zimmer genannt: 2P
    \item Mindestens 4 Farbadjektive korrekt verwendet: 2P
    \item Adjektivangleichung korrekt: 1P
    \item Sprachliche Korrektheit: 1P
\end{itemize}

Beispiellösung: ``Mi casa ideal tiene un salón grande con un sofá rojo y una mesa blanca. La cocina es amarilla. En mi dormitorio hay una cama azul y las paredes son verdes.''

% ============================================================================
% 10. ANLAGEN
% ============================================================================
\section{Anlagen}

\begin{enumerate}
    \item Arbeitsblatt (AB-03c.pdf)
    \item Musterlösung (ML-03c.pdf)
    \item Lehrerhinweise (LH-03c.pdf)
    \item Lernpfad (LP-03c.html)
    \item Farbkarten (10 Stück, laminiert)
    \item Bildkarten Möbel (verschiedene Farben)
\end{enumerate}

% ============================================================================
% 11. DIDAKTIK-SCORE
% ============================================================================
\newpage
\section{Didaktik-Score}

\begin{scorebox}
\textbf{Bewertung der didaktischen Qualität nach 10 Dimensionen}\\
\textbf{Ziel-Score: $\geq$ 9.0 (Premium-Qualität)}

\vspace{1em}
\begin{tabular}{|c|l|c|c|p{4.5cm}|}
\hline
\textbf{\#} & \textbf{Dimension} & \textbf{Gew.} & \textbf{Score} & \textbf{Notiz} \\
\hline
1 & Differenzierung & 15\% & 9/10 & [B] Zuordnung mit Hilfe, [S] selbstständig angleichen, [E] freie Beschreibung \\
\hline
2 & Taxonomie (AFB) & 10\% & 9/10 & AFB I: 35\%, AFB II: 45\%, AFB III: 20\% \\
\hline
3 & Lernziel-Alignment & 10\% & 10/10 & Rahmenplan: Adjektivangleichung, Beschreiben \\
\hline
4 & Aktivierung & 10\% & 10/10 & Farbkarten-Drill, PA-Beschreibung, Memory \\
\hline
5 & Scaffolding & 10\% & 9/10 & Tabelle $\rightarrow$ Übung $\rightarrow$ freie Produktion \\
\hline
6 & Feedback-Qualität & 10\% & 9/10 & Elaboriertes Feedback in LP + AB-Vergleich \\
\hline
7 & Sprachliche Klarheit & 10\% & 10/10 & Altersgerecht, klare Regeln \\
\hline
8 & Motivation \& Relevanz & 10\% & 9/10 & ``Mi casa ideal'' = persönlicher Bezug \\
\hline
9 & Inklusion (UDL) & 10\% & 9/10 & Visuell (Karten) + auditiv + kinästhetisch \\
\hline
10 & Praktikabilität & 5\% & 9/10 & 2x45 Min realistisch, Material verfügbar \\
\hline
\multicolumn{3}{|r|}{\textbf{GESAMT:}} & \textbf{9.3/10} & \textcolor{successgreen}{\textbf{FREIGABE}} \\
\hline
\end{tabular}
\end{scorebox}

\end{document}
